\sectionkicker{Source-backed technical notes}
\subsection*{Inference \& Serving — モデルの外側で進む最適化: Technical Notes}
本欄は記事本文で圧縮した一次資料上の情報を、比較・再検証しやすい形で整理したものである。時系列、確認済みの主張、留意点、一次資料URLを掲載する。
\medskip
\subsection*{Theme at a glance}
\addcontentsline{toc}{subsection}{Theme at a glance}
\begin{center}
\footnotesize
\begin{tabularx}{\linewidth}{@{}p{0.20\linewidth}p{0.10\linewidth}p{0.14\linewidth}X@{}}
\toprule
資料 & 位置づけ & 種別 & 時系列 \\
\midrule
Jalapeño inference chip & 主要資料 & OTHER & 2026-06-24 (HARDWARE\_ANNOUNCEMENT) \\
vLLM v0.24.0 & 主要資料 & FRAMEWORK & 2026-06-29 (PROJECT\_RELEASE) \\
SGLang v0.5.14 & 主要資料 & FRAMEWORK & 2026-06-26 (PROJECT\_RELEASE) \\
FlashInfer v0.6.13 & 主要資料 & FRAMEWORK & 2026-06-24 (PROJECT\_RELEASE) \\
Moebius: Serving Mixture-of-Expert Models with Seamless Runtime Parallelism Switch & 補足資料 & 論文 & 2026-06-25 (論文\_RELEASE) \\
HBM Is Not All You Need: Efficient Disaggregated LLM Serving across Memory-heterogeneous Accelerators & 補足資料 & 論文 & 2026-06-29 (論文\_RELEASE) \\
\bottomrule
\end{tabularx}
\end{center}
\normalsize

\begin{technicalnote}{Jalapeño inference chip}{主要資料}
\begin{tabularx}{\linewidth}{@{}>{\bfseries}p{0.22\linewidth}X@{}}
組織 & OpenAI \& Broadcom \\
種別 & OTHER \\
時系列 & 2026-06-24 (HARDWARE\_ANNOUNCEMENT) \\
\end{tabularx}
\smallskip
{\bfseries 一次資料から整理したtechnical points}
\begin{itemize}[leftmargin=1.5em,itemsep=0.35em]
\item \textbf{Vendor claim}: OpenAI and Broadcom introduced Jalapeño, a custom LLM inference processor; engineering samples were running ML workloads at target frequency/power while final performance measurement was still underway.
\end{itemize}
{\bfseries 読む際の境界}
\begin{itemize}[leftmargin=1.5em,itemsep=0.35em]
\item \textbf{分析上の留意点}: The source is first-party OpenAI material. Capability, benchmark, efficiency, or performance statements remain vendor claims unless explicitly described as externally reproduced.
\end{itemize}
{\bfseries 一次資料}
\begingroup\sloppy
\begin{itemize}[leftmargin=1.5em,itemsep=0.25em]
\item {\scriptsize\url{https://openai.com/index/openai-broadcom-jalapeno-inference-chip}}
\end{itemize}
\endgroup
\end{technicalnote}
\medskip

\begin{technicalnote}{vLLM v0.24.0}{主要資料}
\begin{tabularx}{\linewidth}{@{}>{\bfseries}p{0.22\linewidth}X@{}}
組織 & vLLM Project \\
種別 & FRAMEWORK \\
時系列 & 2026-06-29 (PROJECT\_RELEASE) \\
\end{tabularx}
\smallskip
{\bfseries 一次資料から整理したtechnical points}
\begin{itemize}[leftmargin=1.5em,itemsep=0.35em]
\item \textbf{Project claim}: The release adds MiniMax-M3 support and continues DeepSeek-V4, Model Runner V2, diffusion-model and expert-parallel serving work.
\end{itemize}
{\bfseries 読む際の境界}
\begin{itemize}[leftmargin=1.5em,itemsep=0.35em]
\item \textbf{分析上の留意点}: GitHub release notes are primary project evidence for release contents, but project performance claims are not independently reproduced here.
\end{itemize}
{\bfseries 一次資料}
\begingroup\sloppy
\begin{itemize}[leftmargin=1.5em,itemsep=0.25em]
\item {\scriptsize\url{https://github.com/vllm-project/vllm/releases/tag/v0.24.0}}
\end{itemize}
\endgroup
\end{technicalnote}
\medskip

\begin{technicalnote}{SGLang v0.5.14}{主要資料}
\begin{tabularx}{\linewidth}{@{}>{\bfseries}p{0.22\linewidth}X@{}}
組織 & SGLang Project \\
種別 & FRAMEWORK \\
時系列 & 2026-06-26 (PROJECT\_RELEASE) \\
\end{tabularx}
\smallskip
{\bfseries 一次資料から整理したtechnical points}
\begin{itemize}[leftmargin=1.5em,itemsep=0.35em]
\item \textbf{Project claim}: The release adds support for GLM-5.2, Kimi K2.7 Code and other new models, together with serving/kernel changes including DeepSeek-family work.
\end{itemize}
{\bfseries 読む際の境界}
\begin{itemize}[leftmargin=1.5em,itemsep=0.35em]
\item \textbf{分析上の留意点}: GitHub release notes are primary project evidence for release contents, but project performance claims are not independently reproduced here.
\end{itemize}
{\bfseries 一次資料}
\begingroup\sloppy
\begin{itemize}[leftmargin=1.5em,itemsep=0.25em]
\item {\scriptsize\url{https://github.com/sgl-project/sglang/releases/tag/v0.5.14}}
\end{itemize}
\endgroup
\end{technicalnote}
\medskip

\begin{technicalnote}{FlashInfer v0.6.13}{主要資料}
\begin{tabularx}{\linewidth}{@{}>{\bfseries}p{0.22\linewidth}X@{}}
組織 & FlashInfer Project \\
種別 & FRAMEWORK \\
時系列 & 2026-06-24 (PROJECT\_RELEASE) \\
\end{tabularx}
\smallskip
{\bfseries 一次資料から整理したtechnical points}
\begin{itemize}[leftmargin=1.5em,itemsep=0.35em]
\item \textbf{Project claim}: The release contains Blackwell/GQA decode, attention/MoE, quantization and autotuning kernel/backend changes.
\end{itemize}
{\bfseries 読む際の境界}
\begin{itemize}[leftmargin=1.5em,itemsep=0.35em]
\item \textbf{分析上の留意点}: GitHub release notes are primary project evidence for release contents, but project performance claims are not independently reproduced here.
\end{itemize}
{\bfseries 一次資料}
\begingroup\sloppy
\begin{itemize}[leftmargin=1.5em,itemsep=0.25em]
\item {\scriptsize\url{https://github.com/flashinfer-ai/flashinfer/releases/tag/v0.6.13}}
\end{itemize}
\endgroup
\end{technicalnote}
\medskip

\begin{technicalnote}{Moebius: Serving Mixture-of-Expert Models with Seamless Runtime Parallelism Switch}{補足資料}
\begin{tabularx}{\linewidth}{@{}>{\bfseries}p{0.22\linewidth}X@{}}
組織 & Paper authors \\
種別 & 論文 \\
時系列 & 2026-06-25 (論文\_RELEASE) \\
\end{tabularx}
\smallskip
{\bfseries 一次資料から整理したtechnical points}
\begin{itemize}[leftmargin=1.5em,itemsep=0.35em]
\item \textbf{Author claim}: The authors propose Moebius to switch MoE serving between expert and tensor parallel layouts at runtime without restarting the engine.
\end{itemize}
{\bfseries 読む際の境界}
\begin{itemize}[leftmargin=1.5em,itemsep=0.35em]
\item \textbf{分析上の留意点}: The evidence is the authors' arXiv paper/model card and reported evaluations; results are not independently reproduced in this Evidence review.
\end{itemize}
{\bfseries 一次資料}
\begingroup\sloppy
\begin{itemize}[leftmargin=1.5em,itemsep=0.25em]
\item {\scriptsize\url{http://arxiv.org/abs/2606.26607v1}}
\end{itemize}
\endgroup
\end{technicalnote}
\medskip

\begin{technicalnote}{HBM Is Not All You Need: Efficient Disaggregated LLM Serving across Memory-heterogeneous Accelerators}{補足資料}
\begin{tabularx}{\linewidth}{@{}>{\bfseries}p{0.22\linewidth}X@{}}
組織 & Paper authors \\
種別 & 論文 \\
時系列 & 2026-06-29 (論文\_RELEASE) \\
\end{tabularx}
\smallskip
{\bfseries 一次資料から整理したtechnical points}
\begin{itemize}[leftmargin=1.5em,itemsep=0.35em]
\item \textbf{Author claim}: The authors propose HMA-Serve for disaggregated LLM serving across memory-heterogeneous accelerators.
\end{itemize}
{\bfseries 読む際の境界}
\begin{itemize}[leftmargin=1.5em,itemsep=0.35em]
\item \textbf{分析上の留意点}: The evidence is the authors' arXiv paper/model card and reported evaluations; results are not independently reproduced in this Evidence review.
\end{itemize}
{\bfseries 一次資料}
\begingroup\sloppy
\begin{itemize}[leftmargin=1.5em,itemsep=0.25em]
\item {\scriptsize\url{http://arxiv.org/abs/2606.29986v1}}
\end{itemize}
\endgroup
\end{technicalnote}
\medskip
